\documentclass[10pt,a4paper]{amsart}
\usepackage[margin=19mm]{geometry}
\usepackage{amsmath,amssymb,mathtools}
\usepackage[scr=rsfs,cal=euler]{mathalfa}
\usepackage{xfrac}
\usepackage[unicode,hidelinks,bookmarks=false]{hyperref}
\newcommand{\ie}{i.e.\ }
\newcommand{\cf}{cf.\ }
\newcommand{\resp}{respectively\ }
\newcommand{\Subsection}{\subsection}
\providecommand{\nobreakdash}{\nobreak\hbox{-}}
\setlength{\emergencystretch}{2em}
\multlinegap0pt
\usepackage{mathtools}
\usepackage{amssymb}
\newtheorem{lemma}{Lemma}
\newcommand{\qbinom}[2]{\genfrac{[}{]}{0pt}{}{#1}{#2}_{\!q}} % Gaussian binomial
\title{Algebraic geometric framework of Rogers--Ramanujan identities}
\author{Yifeng Huang, Kenny Lau, Ken Ono, Peter Paule}
\date{Source: arXiv:2608.15219v1 (CC BY 4.0)}
\begin{document}
\maketitle

\noindent\emph{Source and licence.} Algebraic geometric framework of Rogers--Ramanujan identities. Version arXiv:2608.15219v1 (CC BY 4.0), distributed by arXiv under the Creative Commons Attribution 4.0 International licence, http://creativecommons.org/licenses/by/4.0/, as recorded in the arXiv licence metadata for this exact version and shown on the abstract page https://arxiv.org/abs/2608.15219v1. Full licence notice: \texttt{licenses/arxiv-2608.15219-CC-BY-4.0.txt}.\par\smallskip

\noindent\textit{Editorial scope.} Counted unit: the article's lemma lem:q-vander-k (source0.tex lines 482--487). The article's complete original proof of it is delivered with it (source0.tex lines 717--802, one \texttt{proof} environment of 86 lines, which the article places away from the statement). No mathematics is written, repaired or re-derived: the statement and the proof are the article's own lines, in the article's own notation, and no retained line is substituted or re-flowed.  No further source-local context is needed: every cross-reference of every retained line resolves inside the two delivered spans.  Nothing was omitted from any retained span, so the package carries no omission record.  No retained line contains a citation, so the wrapper carries no bibliography and no bibliography entry.  No retained line contains a figure environment, an image-inclusion command or drawing code, so no image is delivered and the package declares no asset dependency.  Editorial formatting only, in addition: an AMS \texttt{amsart} preamble that restates the article's own declarations and macros used by the retained lines, namely the environments \texttt{lemma} on separate counters, together with the standard packages those lines need.  The retained environments print in this document as Lemma 1 in order of appearance, whereas the article prints the same items with the numbering of its own full document.  The complete native source of the article as distributed in the arXiv e-print for this exact version is bundled unchanged as \texttt{sources/207723/source0.tex}.
\begin{lemma}[$k$-part $q$-Vandermonde Identity]\label{lem:q-vander-k}
Let $N = \sum_{i=1}^k N_i$. Then
\begin{equation}\label{eq:q-vander-k}
    \qbinom{N}{R} = \sum_{r_1 + \dots + r_k = R} q^{\sum_{1\le i < j \le k} (N_i - r_i)r_j} \prod_{i=1}^k \qbinom{N_i}{r_i}.
\end{equation}
\end{lemma}

\begin{proof}
Let $S(q)$ denote the left-hand side.
% , and write
% \[
%     \Phi(r,m)=r^2-rm+m^2.
% \]
Applying Lemma~\ref{lem:q-vander-k} to the last $q$-binomial coefficient with block ordering
\(
    (r_1-r_2,c), (r_2-n_1,a-n_1),
\)
and using
\[
    \qbinom{r_2}{n_1}\qbinom{r_2-n_1}{a-n_1}
    =\qbinom{r_2}{a}\qbinom{a}{n_1},
\]
we obtain
\[
    S(q)=\sum_{n_1,n_2,a,c,n_5}q^{Q_1}
    \qbinom{a}{n_1}\qbinom{n_5}{n_2}
    \qbinom{r_1-r_2}{c}
    \qbinom{r_2}{a}\qbinom{r_2}{n_5},
\]
where
\[
    Q_1=
    \left.Q(\mathbf n)\right|_{n_4=a+c}
    +(r_1-r_2-c)(a-n_1).
\]

Set
\[
    m_1=n_1+n_2+c,\qquad m_2=a+n_5.
\]
The $q$-Vandermonde identity with block ordering
\[
    (a,n_1),\qquad(n_5,n_2),\qquad(r_1-r_2,c)
\]
gives
\[
    \sum_{n_1+n_2+c=m_1}q^{E_1}
    \qbinom{a}{n_1}\qbinom{n_5}{n_2}
    \qbinom{r_1-r_2}{c}
    =\qbinom{r_1-r_2+m_2}{m_1},
\]
where
\[
    E_1=(a-n_1)n_2+(a-n_1)c+(n_5-n_2)c.
\]

A direct expansion gives
\[
    Q_2\coloneqq Q_1-E_1=r_1^2-r_1 m_1+m_1^2
    +a^2+an_5+n_5^2-ar_2+r_2^2,
\]
which does not depend on $n_1,n_2,c$.
Consequently,
\begin{equation}\label{eq:b=7}
    S(q)=\sum_{m_1,a,n_5}
    q^{Q_2}
    \qbinom{r_1-r_2+m_2}{m_1}
    \qbinom{r_2}{a}\qbinom{r_2}{n_5}.
\end{equation}

Now put $m_2=a+n_5$. The $q$-Vandermonde identity with block ordering
$ (r_2,a),(r_2,n_5)$ yields
\[
    \sum_{a+n_5=m_2}q^{E_2}
    \qbinom{r_2}{a}\qbinom{r_2}{n_5}
    =\qbinom{2r_2}{m_2},
\]
where $E_2=(r_2-a)n_5$.  Whenever $a+n_5=m_2$, we have
\[
    Q_3\coloneqq Q_2-E_2
    =r_1^2-r_1 m_1+m_1^2+r_2^2-r_2 m_2+m_2^2,
\]
which again does not depend on $a,n_5$. 

Grouping the right-hand side of \eqref{eq:b=7} by $m_2$ therefore gives
\[
    S(q)=\sum_{m_1,m_2}
    q^{Q_3}
    \qbinom{r_1-r_2+m_2}{m_1}
    \qbinom{2r_2}{m_2},
\]
which is the desired identity.
\end{proof}

\end{document}
